\section{Lean formalization architecture}
\label{sec:lean-plan}

\subsection{Integration target}

The ordinary implementation should live under
\begin{center}
\code{Logic/PrimitiveRecursiveArithmetic/Derivations/Lean/}\allowbreak
\code{IPRADerivations/Framework/SCOPE/}.
\end{center}
It should be exposed by an umbrella import only after every module builds.  Foundation adapters belong in \code{IPRABridge} and should reuse the already checked open-sequent transports.

\subsection{Structural kernel}

The archive's \code{lean/SCOPE/StructuralKernel.lean} is a no-placeholder source-level candidate.  Its core declarations have the following shapes:

\begin{lstlisting}
theorem replaceContext (h : Gamma |- A)
    (supply : forall C in Gamma, Delta |- C) : Delta |- A :=
  h.trans supply

structure Contract (n) (Gamma) (goal) where
  valid   : Fm PTm (n + 1)
  extract : valid :: Gamma.map Fm.shift |- goal.shift

theorem Contract.close
    (producer : Gamma |- Fm.ex C.valid) : Gamma |- goal :=
  .exE producer C.extract

structure ResidualResult (n) (Gamma) (goal) where
  residuals : List (Fm PTm n)
  derivation : residuals ++ Gamma |- goal
\end{lstlisting}

The actual file uses the repository notation \code{Gamma |- [PRAAx] A}.  It adds residual closure, weakening, specialization, boundary introduction, and direct sealing.  It contains no \code{sorry} or new axiom.  It was not elaborated in this runtime because the repository workspace was unavailable locally, so it is not reported as checked.

\subsection{Bounded formula module}

The first substantial module should define the fixed syntax:

\begin{lstlisting}
inductive BForm : Nat -> Type
  | eq    (s t : PTm n)
  | bot
  | and   (A B : BForm n)
  | or    (A B : BForm n)
  | imp   (A B : BForm n)
  | allLt (bound : PTm n) (body : BForm (n+1))
  | exLt  (bound : PTm n) (body : BForm (n+1))
\end{lstlisting}

Define \code{denote} in \(\Fm\) and \code{bit} in \(\PTm\).  The key arbitrary-arity fold construction is:

\begin{lstlisting}
def liftFoldBody (body : PTm (n+1)) : PTm (n+2) :=
  Syn.subst (Fin.cases (var 0) (fun j => var j.succ.succ)) body

def allFoldPR (body : PTm (n+1)) : PR (n+1) :=
  recDef (num 1) (band (var 1) (liftFoldBody body))
\end{lstlisting}

This explicitly resolves the parameter-arity problem.  The compiler should return a package:

\begin{lstlisting}
structure Reflection (A : Fm PTm n) where
  bit      : PTm n
  boolean  : [] |- (bit = 0 or bit = 1)
  forward  : [] |- (A -> bit = 1)
  backward : [] |- (bit = 1 -> A)
\end{lstlisting}

The implementation order is Boolean algebra, fold equations, search adapters, fold reflection, then structural recursion on \code{BForm}.  Nested binder tests must substitute non-variable terms through at least two alternating quantifiers.

\subsection{Least accepted adapter}

Wrap the existing search API in the accepted-bit convention.  A useful structure is:

\begin{lstlisting}
structure LeastAccepted (accepted : PTm (n+1)) (bound : PTm n) where
  witness  : PTm n
  foundBit : PTm n
  boolean  : ...
  sound    : ...
  minimal  : ...
  sentinel : ...
\end{lstlisting}

The implementation should reuse \code{Search.search_found}, \code{search_minimal}, and \code{search_spec}; it should not reprove least search from scratch.

\subsection{Codes and demand closure}

Select one vector representation and prove it before implementing closure.  Mixed radix is attractive because total capacity and lexicographic code order are explicit.  Required modules are:

\begin{enumerate}
\item \code{MixedRadix}: digit, encode, decode, range, and round trip;
\item \code{CappedVector}: coordinate order, join, rank, and cap sum;
\item \code{DemandClosure}: inflationary iteration, strict rank, fixedness, and leastness.
\end{enumerate}

The finite dimension is a Lean parameter or a fixed structure index; object iteration is on the natural step counter.  Coordinatewise boundedness is represented by a reflected fold bit so the induction formula remains QF.

\subsection{Trace module}

A reusable trace package should accept fixed terms/symbols and local derivations rather than a semantic Lean function alone.  One possible interface is:

\begin{lstlisting}
structure TraceContract (n : Nat) where
  statePR  : PR (n+1)
  init     : PTm n
  inv      : Fm PTm (n+1)
  invQF    : inv.QF
  base     : Gamma |- inv[init]
  step     : inv :: Gamma.shift |- inv[next]
\end{lstlisting}

The precise indexing should be chosen to minimize rewrites with \code{recDef}.  The result theorem applies \code{IPRA.ind} and then \code{Prf.subst} at arbitrary fuel.

\subsection{Finite strategies and viability}

Begin with the \(\forall\exists\) case.  Encode a table by mixed radix and prove table lookup equations.  Then add a host-level datatype for fixed quantifier shapes and compile it recursively.

Viability should be a separate module.  Its acceptance test must include the dead-end example from \cref{sec:viability}.  Soundness and completeness should be stated as ordinary object derivations of reflected path predicates, not merely as Lean propositions over lists.

\subsection{Observers and scope effects}

Define structures for:

\begin{itemize}
\item exact observer to a fixed outer term;
\item explicit congruence observer;
\item rational metric diameter observer;
\item gap-guarded finite label observer.
\end{itemize}

The scope compiler can use a host enum
\begin{lstlisting}
inductive Escape
  | private | existential | stable | canonical | residual
\end{lstlisting}
but each constructor must require the proof payload corresponding to its semantics.  In particular, \code{stable} is not a Boolean annotation; it carries an observer derivation.

\subsection{Graph compiler}

The graph compiler is a Lean metaprogram or ordinary recursive function over a fixed typed graph.  Every node returns a term, formula, and derivation package.  The final result is a \code{ResidualResult}.  The compiler may fail with a structured diagnostic when a resource cap, observer proof, or producer is missing.

A useful diagnostic is the minimal residual frontier: traverse backward from the goal, remove producers discharged internally, and report only reachable unsupplied leaves.  This makes failed compilation mathematically informative.

\subsection{Acceptance gates}

The formalization plan in the archive lists the complete gates.  The most important are:

\begin{itemize}
\item focused and full \code{lake build} success;
\item no \code{sorry}, \code{admit}, new \code{axiom}, or unsafe Boolean reflection;
\item \code{#print axioms} audit consistent with repository policy;
\item substitution at compound terms;
\item nested bounded alternation;
\item a cyclic demand example requiring multiple passes;
\item a greedy-dead-end viability example;
\item two distinct valid certificates with one stable public observation;
\item an exact residual visible in the final theorem type;
\item replay through both scoped bridge directions.
\end{itemize}

\subsection{Four staged pilots}

After the core builds, migrate:

\begin{enumerate}
\item a finite arithmetic theorem with a bounded certificate;
\item the quantitative bisection contract of \cref{sec:trace};
\item a finite complex or cover example using tree descent;
\item a height-bounded Diophantine list with the height theorem left residual.
\end{enumerate}

These pilots exercise every layer without requiring an immediate attack on the hardest analytic or Diophantine producers.
